\documentclass{exam}
\usepackage{graphicx}
\usepackage[utf8]{inputenc}
\usepackage[main=english,vietnamese]{babel}
\usepackage{amsmath}
\usepackage{hyperref}
\usepackage{amsthm}
\usepackage{tcolorbox}
\usepackage{amsfonts}
\usepackage{amssymb}
\usepackage{mathrsfs}
\usepackage{centernot}
\usepackage{cases}
\usepackage{physics}
\usepackage[shortlabels]{enumitem}
\usepackage{tikz}
\usepackage{lipsum}
\usepackage{sidecap}
\usepackage{verbatim}
\usepackage{listings}
\usepackage{kvmap}

\usetikzlibrary{calc}
\usetikzlibrary{decorations.pathmorphing}

\definecolor{borderblue}{HTML}{33c2ff}
\definecolor{codegreen}{rgb}{0,0.6,0}
\definecolor{codegray}{rgb}{0.5,0.5,0.5}
\definecolor{codepurple}{rgb}{0.58,0,0.82}
\definecolor{backcolour}{rgb}{0.95,0.95,0.92}

\lstdefinestyle{mystyle}{
    backgroundcolor=\color{backcolour},   
    commentstyle=\color{codegreen},
    keywordstyle=\color{magenta},
    numberstyle=\tiny\color{codegray},
    stringstyle=\color{codepurple},
    basicstyle=\ttfamily\footnotesize,
    breakatwhitespace=false,         
    breaklines=true,                 
    captionpos=b,                    
    keepspaces=true,                 
    numbers=left,                    
    numbersep=5pt,                  
    showspaces=false,                
    showstringspaces=false,
    showtabs=false,
    tabsize=2
}

\lstset{style=mystyle}

\newcommand{\paren}[1]{\left(#1\right)}
\newcommand{\curly}[1]{\left\{#1\right\}}

\allowdisplaybreaks

\makeatletter
\long\def\paragraph{%
  \@startsection{paragraph}{4}%
  {\z@}{2ex \@plus 1ex \@minus .2ex}{-1em}%
  {\normalfont\normalsize\bfseries}%
}
\makeatother

\makeatletter
\def\@xequals@fill{\arrowfill@\Relbar\Relbar\Relbar}
\newcommand*\xequals[2][]{\DOTSB\ext@arrow0055\@xequals@fill{#1}{#2}}
\makeatother

\makeatletter
\renewcommand*\env@matrix[1][*\c@MaxMatrixCols c]{%
  \hskip -\arraycolsep
  \let\@ifnextchar\new@ifnextchar
  \array{#1}}
\makeatother

\NewTColorBox{subcircuit}{m}{
  standard jigsaw,
  sharp corners,
  boxrule=0.4pt,
  coltitle=black,
  colframe=black,
  opacityback=0,
  opacitybacktitle=0,
  fonttitle=\normalfont\bfseries\upshape,
  fontupper=\normalfont,
  title={Subcircuit \texttt{#1}},
  after title={.},	
  attach title to upper={\ },
}

\newcommand{\lnand}{\uparrow}
\newcommand{\lnor}{\downarrow}
\newcommand{\lxor}{\oplus}

\title{ECE341 - Homework 7}
\author{Ton-Minh-Ky Tran}
\date{\today}

\begin{document}

\maketitle
\section{Unconditional jump instruction}
The \verb|J| instruction corresponds to an opcode of \verb|000010|. By inserting an instruction with the opcode and $2$ as the label corresponding to an absolute jump to \verb|0x8|, the yielded program jumps to . The program is found in the ROM image \verb|loop.txt|.

\section{While-loop using branches and jumps}
We simply look up the MIPS instruction encoding to encode each given instruction. The jump address is \verb|0x3|. The program is found in the ROM image \verb|whileloop.txt|.

\section{Null-terminated string storage in RAM}
The starting address is \verb|$r2 := 0x30| and we wish to move each byte of the string (two characters) to \verb|r2 + 2i|, where \verb|i| is the displacement of the starting address signifying number of bytes stored. Using \verb|r1| to store two characters at a time and then move it to the address stored in \verb|r2|, we yield what we desire---the string \verb|Hello: Ky\0| in the address \verb|0x30|. The program is found in the ROM image \verb|hello.txt|.

\section{Load null-terminated string in RAM to TTY console}
Using the same procedure outlined in \textbf{Exercise 3} to preload the string into RAM, we know that the string address \verb|S| is \verb|0x30|, the string length \verb|n| is \verb|0xa|, and the memory-mapped address of the console \verb|D| is \verb|0x80000008|. From this point, this is just a loop that reads the character in RAM using the \verb|lb| instruction to store in a register and then load the character to \verb|D|. The loop terminates if \verb|P| reaches \verb|S + n|, after which we load the line feed character to \verb|D|. The program is found in the ROM image \verb|println.txt|.

\end{document}